\documentclass[11pt]{article}

\usepackage[T1]{fontenc}
\usepackage[utf8]{inputenc}
\usepackage{lmodern}
\usepackage[margin=1in]{geometry}
\usepackage{microtype}
\usepackage{amsmath,amssymb,mathtools}
\usepackage{enumitem}
\usepackage[hidelinks]{hyperref}

\newcommand{\E}{\mathbb{E}}
\newcommand{\Var}{\operatorname{Var}}
\newcommand{\Cov}{\operatorname{Cov}}
\newcommand{\USD}{\text{USD}}

\title{A Task-Level Provider Utility Model}
\author{Matthew Pollard}
\date{}

\begin{document}
\maketitle

Provider selection should minimize the expected risk-adjusted cost of completing an
entire task, not the nominal price of an isolated request. A task consists of $N$
conversational turns whose reusable prefix generally grows over time. The model
includes cached and uncached input charges, output-token charges, latency and
token-generation time converted into dollars, endpoint failures and retries,
stochastic cache misses, shrinkage of noisy endpoint statistics, and a variance
penalty for operational and estimation risk.

\section{Conversation structure}

For turn $k=1,\ldots,N$, let
\begin{itemize}[nosep]
  \item $A_k$ be newly appended prompt tokens, such as a user message or tool result;
  \item $O_k$ be completion tokens generated on the turn; and
  \item $S_k$ be reusable prefix tokens preceding the new input.
\end{itemize}

For a conversation that resubmits its complete history,
\begin{equation}
S_k=\sum_{j<k}(A_j+O_j).
\end{equation}
The complete prompt length is therefore
\begin{equation}
P_k=S_k+A_k.
\end{equation}
This distinction matters: $A_k$ is necessarily new input, whereas only $S_k$
is eligible for a cache hit.

\section{Endpoint parameters}

For endpoint $e$, let
\begin{itemize}[nosep]
  \item $w_e$ be the uncached-input or cache-write price per token;
  \item $r_e$ be the cached-read price per token;
  \item $c_e$ be the completion price per token;
  \item $f_e$ be the fixed request charge;
  \item $h_e$ be the probability that the reusable prefix is cached;
  \item $p_e=1-h_e$ be the cache-miss probability;
  \item $u_e$ be the probability that an individual attempt succeeds;
  \item $L_e$ be time to first token in seconds;
  \item $G_e$ be token-generation throughput in tokens per second;
  \item $T_e^F$ be time consumed by a failed attempt;
  \item $\beta_e\in[0,1]$ be the fraction of attempted monetary input cost charged
        on failure; and
  \item $v$ be the user's value of time in USD per hour.
\end{itemize}

Define the dollar value of one second as
\begin{equation}
\rho=\frac{v}{3600}.
\end{equation}
The successful-turn wall-clock time is
\begin{equation}
T_{e,k}^S=L_e+\frac{O_k}{G_e}.
\end{equation}
The first term values latency to first token; the second values token-generation
time.

\section{Cache-miss model}

For every attempt $j$ at turn $k$, define
\begin{equation}
M_{k,j}\sim\operatorname{Bernoulli}(p_e),
\end{equation}
where $M_{k,j}=1$ denotes a complete miss of the reusable prefix. Cache events
are assumed independent across attempts and turns and independent of endpoint
availability.

Let
\begin{equation}
\Delta_e=w_e-r_e
\end{equation}
be the cache-miss price spread. The attempted input cost is
\begin{equation}
I_{e,k,j}
=f_e+w_eA_k+S_k\left[r_e+M_{k,j}\Delta_e\right].
\end{equation}
Thus new tokens $A_k$ are always charged at $w_e$. On a hit, the reusable
prefix $S_k$ is charged at $r_e$; on a miss, the same prefix is charged at
$w_e$.

For a fixed hit probability,
\begin{align}
\bar I_{e,k}
&=\E[I_{e,k,j}] \\
&=f_e+w_eA_k+S_k\left[r_e+p_e\Delta_e\right],
\end{align}
and
\begin{equation}
V_{e,k}^I
=\Var(I_{e,k,j})
=S_k^2\Delta_e^2p_e(1-p_e).
\end{equation}

\section{Uptime and retry model}

Let an attempt succeed independently with probability $u_e$. The number of
failed attempts before the first success is
\begin{equation}
K_{e,k}\sim\operatorname{GeometricFailures}(u_e),
\end{equation}
with
\begin{equation}
\E[K_{e,k}]=\frac{1-u_e}{u_e},
\qquad
\Var(K_{e,k})=\frac{1-u_e}{u_e^2}.
\end{equation}
This models retrying the endpoint until one attempt succeeds. An implementation
may replace this with an explicit fallback chain when the fallback endpoint is
known.

The cost of the successful attempt is
\begin{equation}
Y_{e,k}^S=I_{e,k}+c_eO_k+\rho T_{e,k}^S.
\end{equation}
The cost of a failed attempt is
\begin{equation}
Y_{e,k}^F=\beta_e I_{e,k}+\rho T_e^F.
\end{equation}
The parameter $\beta_e$ makes the billing assumption explicit. Setting
$\beta_e=0$ means failed attempts incur time cost but no token charge, whereas
$\beta_e=1$ conservatively assigns them their full attempted input charge.

The completed-turn cost is
\begin{equation}
X_{e,k}=Y_{e,k}^S+\sum_{j=1}^{K_{e,k}}Y_{e,k,j}^F.
\end{equation}
Define
\begin{equation}
B_e=\frac{1-u_e}{u_e}.
\end{equation}
Then
\begin{equation}
\mu_{e,k}^S
=\bar I_{e,k}+c_eO_k
+\rho\left(L_e+\frac{O_k}{G_e}\right),
\end{equation}
\begin{equation}
\mu_{e,k}^F=\beta_e\bar I_{e,k}+\rho T_e^F,
\end{equation}
and the expected cost of completing turn $k$ is
\begin{equation}
\boxed{
\E[X_{e,k}]=\mu_{e,k}^S+B_e\mu_{e,k}^F
}.
\end{equation}
Equivalently,
\begin{equation}
\boxed{
\begin{aligned}
\E[X_{e,k}]
={}&\bar I_{e,k}+c_eO_k
+\rho\left(L_e+\frac{O_k}{G_e}\right) \\
&+\frac{1-u_e}{u_e}
\left(\beta_e\bar I_{e,k}+\rho T_e^F\right).
\end{aligned}
}
\end{equation}
Unlike a simple $(1-u_e)$ surcharge, the factor $(1-u_e)/u_e$ accounts for
the possibility of more than one failed attempt.

\section{Task-level expected cost}

The task cost is the sum of its completed-turn costs:
\begin{equation}
X_e^{\mathrm{task}}=\sum_{k=1}^{N}X_{e,k}.
\end{equation}
Therefore,
\begin{equation}
\boxed{
\mu_e^{\mathrm{task}}
=\E[X_e^{\mathrm{task}}]
=\sum_{k=1}^{N}\left[\mu_{e,k}^S+B_e\mu_{e,k}^F\right]
}.
\end{equation}
This expected dollar cost per completed task should be the primary ranking
quantity.

\section{Constant-growth conversation}

For intuition, suppose every turn adds $a$ new prompt tokens and produces $o$
completion tokens. Let
\begin{equation}
d=a+o,
\qquad
S_k=(k-1)d.
\end{equation}
Define
\begin{equation}
\pi_e
=r_e+p_e(w_e-r_e)
=h_er_e+(1-h_e)w_e,
\end{equation}
\begin{equation}
b_e=f_e+w_ea,
\end{equation}
\begin{equation}
y_e=c_eo+\rho\left(L_e+\frac{o}{G_e}\right),
\end{equation}
\begin{equation}
z_e=\rho T_e^F,
\end{equation}
and
\begin{equation}
D_e=1+\beta_e\frac{1-u_e}{u_e}.
\end{equation}
Then the expected task cost has the closed form
\begin{equation}
\boxed{
\mu_e^{\mathrm{task}}
=N\left[D_eb_e+y_e+B_ez_e\right]
+D_e\pi_ed\frac{N(N-1)}{2}
}.
\end{equation}
The first term is linear in the number of turns. The second is quadratic because
the prefix grows while being submitted repeatedly.

The expected cache-miss premium, excluding the ordinary cached-read charge, is
\begin{equation}
\boxed{
D_ep_e(w_e-r_e)d\frac{N(N-1)}{2}
}.
\end{equation}
Consequently, a 5\% cache-miss rate is not merely a 5\% task-level surcharge. It
applies to a sequence of progressively larger prefixes.

\section{Operational variance}

Assume successful and failed attempt times have variances
\begin{equation}
\Var(T_{e,k}^S)=(\sigma_{e,k}^S)^2,
\qquad
\Var(T_e^F)=(\sigma_e^F)^2.
\end{equation}
Then
\begin{equation}
V_{e,k}^S=V_{e,k}^I+\rho^2(\sigma_{e,k}^S)^2,
\end{equation}
and
\begin{equation}
V_{e,k}^F=\beta_e^2V_{e,k}^I+\rho^2(\sigma_e^F)^2.
\end{equation}
Using the variance of a random sum,
\begin{equation}
\boxed{
\Var(X_{e,k})
=V_{e,k}^S+B_eV_{e,k}^F
+\frac{1-u_e}{u_e^2}\left(\mu_{e,k}^F\right)^2
}.
\end{equation}
The three components are uncertainty in the successful attempt, uncertainty in
the costs of failed attempts, and uncertainty in how many failures occur before
success.

Under independence across turns,
\begin{equation}
\boxed{
V_{e,\mathrm{process}}^{\mathrm{task}}
=\sum_{k=1}^{N}\Var(X_{e,k})
}.
\end{equation}
For constant transcript growth,
\begin{equation}
\sum_{k=1}^{N}S_k^2
=d^2\frac{N(N-1)(2N-1)}{6}.
\end{equation}
Accordingly, cache-miss variance grows as $O(N^3)$, and its standard deviation
grows as $O(N^{3/2})$.

\section{Shrinkage of cache hit rate and uptime}

Raw 24-hour endpoint statistics can be unstable for low-volume endpoints. Both
cache hit rate and uptime should therefore be shrunk toward population priors.

For cache hits, let $\hat h_e$ be the published cache-hit rate, $n_e^h$ the
effective cache sample size, $h_0$ the population prior mean, and $\kappa_h$
the prior strength. Use the beta prior
\begin{equation}
h_e\sim\operatorname{Beta}
\left(\kappa_hh_0,\,\kappa_h(1-h_0)\right).
\end{equation}
After observing effective evidence $n_e^h$, the posterior parameters are
\begin{equation}
\alpha_e^h=\kappa_hh_0+n_e^h\hat h_e,
\end{equation}
and
\begin{equation}
\beta_e^h=\kappa_h(1-h_0)+n_e^h(1-\hat h_e).
\end{equation}
The shrunk hit rate is
\begin{equation}
\boxed{
\tilde h_e
=\frac{\kappa_hh_0+n_e^h\hat h_e}{\kappa_h+n_e^h}
}.
\end{equation}
Its posterior variance is
\begin{equation}
\boxed{
s_{h,e}^2
=\frac{\alpha_e^h\beta_e^h}
{(\alpha_e^h+\beta_e^h)^2(\alpha_e^h+\beta_e^h+1)}
}.
\end{equation}

The effective sample size should not simply equal raw token volume. Cache outcomes
within one request or conversation are correlated. A practical approximation is
\begin{equation}
n_e^h=\frac{T_e}{B_h},
\end{equation}
where $T_e$ is relevant prompt-token volume and $B_h$ is a configurable
correlation-block size.

Uptime is treated analogously. Given published uptime $\hat u_e$, effective
sample size $n_e^u$, prior mean $u_0$, and prior strength $\kappa_u$,
\begin{equation}
\boxed{
\tilde u_e
=\frac{\kappa_uu_0+n_e^u\hat u_e}{\kappa_u+n_e^u}
}.
\end{equation}
Its posterior variance $s_{u,e}^2$ follows the same beta-posterior expression.
Uptime evidence should be measured in requests or approximately independent time
buckets, not tokens.

The expected-cost and process-variance equations use the shrunk values
\begin{equation}
h_e=\tilde h_e,
\qquad
u_e=\tilde u_e.
\end{equation}

\section{Shrinkage of latency and generation speed}

Latency and generation speed can also be unstable at low volume. Because both
quantities are positive and usually right-skewed, shrink them in log space. For a
positive performance parameter $\theta_e$, such as TTFT or seconds per generated
token, let $\hat\theta_e$ be the endpoint estimate and $\theta_0$ the model-wide
or provider-wide prior. Define
\begin{equation}
\log\tilde\theta_e
=\omega_e\log\hat\theta_e
+(1-\omega_e)\log\theta_0,
\end{equation}
where
\begin{equation}
\omega_e=\frac{n_e^\theta}{n_e^\theta+\kappa_\theta}.
\end{equation}
For throughput, shrink seconds per token $1/G_e$, rather than throughput $G_e$,
because time cost is linear in $1/G_e$.

When only p50 values are available, the resulting time cost is a point estimate.
When multiple quantiles or raw observations are available, they should also be
used to estimate $\sigma^S$ and $\sigma^F$.

\section{Parameter-estimation risk}

Process variance assumes that the endpoint parameters are known. Shrinkage reduces
estimation error but does not eliminate it. Let
\begin{equation}
\boldsymbol\theta_e=(h_e,u_e,L_e,G_e,T_e^F,\ldots)
\end{equation}
have posterior covariance matrix $\Sigma_{\theta,e}$. A first-order approximation
to uncertainty in expected task cost is
\begin{equation}
\boxed{
V_{e,\mathrm{parameter}}^{\mathrm{task}}
\approx
\nabla_{\!\theta}\mu_e^{\mathrm{task}\,\mathsf T}
\Sigma_{\theta,e}
\nabla_{\!\theta}\mu_e^{\mathrm{task}}
}.
\end{equation}
For independent parameter estimates, this reduces to a sum of squared sensitivities
multiplied by their variances.

Useful sensitivities include
\begin{equation}
\frac{\partial\mu_e^{\mathrm{task}}}{\partial h_e}
=-
\left(1+\beta_e\frac{1-u_e}{u_e}\right)
(w_e-r_e)\sum_{k=1}^{N}S_k,
\end{equation}
\begin{equation}
\frac{\partial\mu_e^{\mathrm{task}}}{\partial u_e}
=-\frac{1}{u_e^2}\sum_{k=1}^{N}\mu_{e,k}^F,
\end{equation}
\begin{equation}
\frac{\partial\mu_e^{\mathrm{task}}}{\partial L_e}=\rho N,
\end{equation}
and, when throughput is treated as a common endpoint parameter,
\begin{equation}
\frac{\partial\mu_e^{\mathrm{task}}}{\partial G_e}
=-\rho\frac{\sum_{k=1}^{N}O_k}{G_e^2}.
\end{equation}

For constant transcript growth,
\begin{equation}
\sum_{k=1}^{N}S_k=O(N^2).
\end{equation}
Therefore uncertainty in the endpoint's true cache-hit rate contributes
\begin{equation}
V_{h,e}^{\mathrm{parameter}}=O(N^4),
\end{equation}
even though independent realized cache misses contribute only $O(N^3)$ process
variance. Equivalently, realized miss-risk standard deviation grows as
$O(N^{3/2})$, while uncertainty about the true miss rate can produce an
$O(N^2)$ standard-deviation penalty. This is why uncertain cache statistics
become especially consequential for long agent trajectories.

\section{Risk-adjusted provider score}

The final provider score is
\begin{equation}
\boxed{
J_e
=\mu_e^{\mathrm{task}}
+\lambda_{\mathrm{process}}
 \sqrt{V_{e,\mathrm{process}}^{\mathrm{task}}}
+\lambda_{\mathrm{parameter}}
 \sqrt{V_{e,\mathrm{parameter}}^{\mathrm{task}}}
}.
\end{equation}
Here $\lambda_{\mathrm{process}}$ penalizes unpredictable realized costs and
delays, while $\lambda_{\mathrm{parameter}}$ penalizes poorly estimated endpoint
statistics. Reasonable interpretations are
\begin{itemize}[nosep]
  \item $\lambda_{\mathrm{process}}=\lambda_{\mathrm{parameter}}=0$: risk-neutral
        expected cost;
  \item $\lambda=1$: a one-standard-deviation penalty; and
  \item $\lambda\approx1.645$: approximate one-sided 95\% protection under a
        normal approximation.
\end{itemize}
Separating the penalties is preferable because a high-volume endpoint can have
meaningful operational variance but little estimation uncertainty, while a
low-volume endpoint may look excellent solely because its statistics are poorly
estimated.

\section{Reported outputs}

For each endpoint, the tool should report
\begin{enumerate}[nosep]
  \item expected dollar cost per completed task;
  \item risk-adjusted dollar score $J_e$;
  \item expected number of failed attempts;
  \item expected uncached-input, cached-read, and output tokens;
  \item dollar cost of TTFT and token generation;
  \item cache-miss premium;
  \item uptime-retry premium;
  \item process standard deviation;
  \item parameter-estimation standard deviation; and
  \item raw and shrunk cache-hit and uptime estimates.
\end{enumerate}

The primary ranking should use dollars per completed task. A normalized
per-million-token figure may be shown secondarily, but it must state whether the
denominator counts unique task tokens or repeatedly submitted prompt tokens.

\end{document}
